\chapter{童年的诗}

玩秋千 石人山情节

2016.9.11

茅草小屋金銮殿，桃杏梨花来装点。 春日争相露笑脸，葛花树上玩秋千。 若有同龄孩童伴，情景定会不一般。 良辰美景一人享，亦开心来亦孤单。

石人山景 石人山景美若画，儿时记忆难忘下。 高山顶峰衬其峻，云雾缭绕摹其神。 谷满林竹花锦绣，潺潺流水汇清泉。

恋家 少时寄住石人山，犹似修隐仙景间。 天亮登山陪桑蚕，入夜土炕头边玩。 山风怒吼夜狼嚎，柴房家犬牙齿露。 一日三餐尚果腹，思乡遥望东方愁。

四伯 追思堂伯父 幼时父母皆亡殁，凄凉处境人怜惜。 尝尽人间冷与暖，变得木讷而寡言。 不养鸡鸭不养猪，不种蔬菜只种谷。 厨中并无盐油醋，舌尖哪有美味留。

无题 民国三十年，恰逢吾降生。 上天不下雨，庄家未收成。 家园遭天灾，倭寇来侵略。 人在襁褓中，母亲患大病。 断了生命水，差点送性命。 - 注：生命之水指母亲奶水。

思往昔看今朝

2015.3

忆往昔峥嵘岁月稠， 想痛处，满脸沟壑老泪流。 看今朝国盛家富， 家和睦儿孙老人事事顺溜。

忆石人涧茅草小屋“金銮殿”

2016.9.9

石人身高八丈三，一半身在云雾间。 石婆闲来无去处，身贴悬崖避风寒。 谷底花树高又俊，秀拔挺立直冲天。 房后园中绿竹林，粗若碗口真稀罕。 门前小河水不断，源自上游山泉涧。 春风吹佛小山村，桃杏梨花齐斗艳。 吾辈此时何处去，葛藤树上荡秋千。

咏春 石人山旧居 茅草小屋金銮殿，桃杏梨花来装点。 三亩翠竹作盆景，门前溪水流不断。 欲问清水何处来，遥指山谷泉水涧。 春色美景独享有，葛花树上荡秋千。

大竹园变堰潭 旧时竹园改了观，现已变成大堰潭。 鱼儿为觅漂浮时，头儿伸出乱眨眼。 若想尝下新鲜味，一网下去捞两篮。

\begin{itemize}[leftmargin=2.5em,itemsep=2pt]
  \item 注：原来的大竹园在58年时，竹子被砍伐，修成了大堰潭，原生态的景观也被破坏殆尽
\end{itemize}

